%-------------------------
% Resume in LaTeX
% Based on: https://github.com/sb2nov/resume (MIT License)
%------------------------

\documentclass[letterpaper,11pt]{article}

\usepackage{latexsym}
\usepackage[empty]{fullpage}
\usepackage{titlesec}
\usepackage{marvosym}
\usepackage[usenames,dvipsnames]{color}
\usepackage{verbatim}
\usepackage{enumitem}
\usepackage[hidelinks]{hyperref}
\usepackage{fancyhdr}
\usepackage[english]{babel}
\usepackage{tabularx}
\input{glyphtounicode}

\pagestyle{fancy}
\fancyhf{}
\fancyfoot{}
\renewcommand{\headrulewidth}{0pt}
\renewcommand{\footrulewidth}{0pt}

\addtolength{\oddsidemargin}{-0.5in}
\addtolength{\evensidemargin}{-0.5in}
\addtolength{\textwidth}{1in}
\addtolength{\topmargin}{-.5in}
\addtolength{\textheight}{1.0in}

\urlstyle{same}
\raggedbottom
\raggedright
\setlength{\tabcolsep}{0in}

\titleformat{\section}{
  \vspace{-4pt}\scshape\raggedright\large
}{}{0em}{}[\color{black}\titlerule \vspace{-5pt}]

\pdfgentounicode=1

\definecolor{resumeblue}{RGB}{70, 130, 180}

\newcommand{\projectLink}[2]{\href{https://moonlitjaderabbit.github.io/#1}{\textbf{#2}}}

\newcommand{\resumeItem}[1]{
  \item\small{
    {#1 \vspace{-2pt}}
  }
}

\newcommand{\resumeSubheading}[4]{
  \vspace{-2pt}\item
    \begin{tabular*}{0.97\textwidth}[t]{l@{\extracolsep{\fill}}r}
      \textbf{#1} & #2 \\
      \textit{\small#3} & \textit{\small #4} \\
    \end{tabular*}\vspace{-7pt}
}

\newcommand{\resumeProjectHeading}[2]{
    \item
    \begin{tabular*}{0.97\textwidth}[t]{l@{\extracolsep{\fill}}r}
      \small#1 & #2 \\
    \end{tabular*}\vspace{-7pt}
}

\renewcommand\labelitemii{$\vcenter{\hbox{\tiny$\bullet$}}$}

\newcommand{\resumeSubHeadingListStart}{\begin{itemize}[leftmargin=0.15in, label={}]}
\newcommand{\resumeSubHeadingListEnd}{\end{itemize}}
\newcommand{\resumeItemListStart}{\begin{itemize}}
\newcommand{\resumeItemListEnd}{\end{itemize}\vspace{-5pt}}

\begin{document}

\begin{center}
    \textbf{\Huge \scshape Junye Zhou} \\ \vspace{4pt}
    \small
    \href{https://moonlitjaderabbit.github.io/}{\color{resumeblue}\textbf{Portfolio}: moonlitjaderabbit.github.io}
    \color{black} $|$
    \href{mailto:junye.zhou@mail.utoronto.ca}{\color{resumeblue}junye.zhou@mail.utoronto.ca}
    \color{black}
\end{center}

\section{Education}
  \resumeSubHeadingListStart
    \resumeSubheading
      {University of Toronto}{Toronto, ON}
      {Bachelor of Applied Science in Computer Engineering}{Expected Graduation: 2030 (PEY Co-op)}
      \resumeItemListStart
        \resumeItem{Coursework: Algorithms \& Data Structures (Python, C), Programming Fundamentals (C++), Digital Systems (Verilog), Electrical Circuits, Intro to Electronics}
      \resumeItemListEnd
  \resumeSubHeadingListEnd

\section{Engineering Experience}
  \resumeSubHeadingListStart
    \resumeSubheading
      {Drone Team Member}{Sep 2026 -- Present}
      {RSX Robotics for Space Exploration}{}
      \resumeItemListStart
        \resumeItem{Developed a 6\,V buck-converter design for the gripper servo and a 19\,V buck-boost converter design to power the Jetson Nano. Prepared the schematics in KiCad}
        \resumeItem{Review battery voltage limits and the power needs of onboard components to guide regulator selection and power distribution planning}
        \resumeItem{Performed soldering work with the electrical subteam. Review UART connections and SiK/MAVLink telemetry for the Kakute H7 flight controller to support ongoing communications work}
      \resumeItemListEnd

    \resumeSubheading
      {Team Lead}{Oct 2026}
      {Trench--Hydro One Innovation Sprint $|$ Winning Concept, 5th Place}{}
      \resumeItemListStart
        \resumeItem{Led a 4-person team that proposed using electromagnetic and acoustic sensors to detect electrical discharges inside sealed capacitor voltage transformers}
        \resumeItem{Used Python and openEMS, with support from AI coding tools, to simulate 300--700\,MHz electromagnetic signals through a simplified oil-filled model and examine the predicted fields at its outer housing}
        \resumeItem{Evaluated detection thresholds in synthetic noise and interference tests; a higher threshold reduced false candidates from 8.93 to 1.02 per record, clarifying the trade-off between false alarms and missed pulses}
      \resumeItemListEnd

    \resumeSubheading
      {Praxis Design Team}{Feb 2026 -- Apr 2026}
      {University of Toronto Praxis II}{}
      \resumeItemListStart
        \resumeItem{Designed a swim-training concept that uses a rotating drum to control rope resistance, targeting 8 levels (60--170\,N) for the Mississauga Aquatic Club}
        \resumeItem{Compared designs for adjustability, ease of manufacture, and poolside safety using Pugh matrices. Wrote requirements and test plans, then presented the team's concept at the Praxis II Showcase}
      \resumeItemListEnd

    \resumeSubheading
      {Bridge Design Team}{Nov 2025 -- Jan 2026}
      {University of Toronto CIV102 $|$ Winning Design, 2nd Place}{}
      \resumeItemListStart
        \resumeItem{Designed, analyzed, and built a 1.25\,m matboard bridge with teammates, placing 2nd in the EngSci class}
        \resumeItem{Checked the team's Python structural analysis against hand calculations to assess the loads the bridge could carry and how it might fail}
      \resumeItemListEnd
  \resumeSubHeadingListEnd

\section{Selected Project}
  \resumeSubHeadingListStart
    \resumeSubheading
      {\href{https://torontodropin.ca/}{Toronto Badminton Drop-in Finder Website}}{University of Toronto}
      {\href{https://torontodropin.ca/}{torontodropin.ca}}{}
      \resumeItemListStart
        \resumeItem{Developed a website using AI coding tools and City of Toronto open data to help badminton players find nearby community centres and track drop-in sessions, serving 70 daily users}
      \resumeItemListEnd
  \resumeSubHeadingListEnd

\section{Technical Skills}
 \begin{itemize}[leftmargin=0.15in, label={}]
    \small{\item{
     \begin{tabularx}{0.97\textwidth}{@{}>{\bfseries}l @{\hspace{1.5em}} X@{}}
      Programming: & C, C++, Python, Verilog, Java \\
      Development: & Git/GitHub, GCC (MSYS2), VS Code, ModelSim, DESim, Quartus Prime \\
      Hardware \& Analysis: & KiCad, LTspice, TI WEBENCH, MATLAB, openEMS, soldering, circuit analysis \\
     \end{tabularx}
    }}
 \end{itemize}

\vspace{0.6em}
\noindent\small For project details and additional work, visit my portfolio:
\href{https://moonlitjaderabbit.github.io/}{https://moonlitjaderabbit.github.io/}

\end{document}
